\documentclass[12pt]{article}
\usepackage{ficha}

\begin{document}

\ficha{Equações Trigonométricas e Fórmulas de Adição}

\instrucoes{Resolve em $\R$, salvo indicação em contrário. Apresenta as soluções gerais com $k \in \Z$.}

\secao{Equações Elementares}{Resolve em $\R$.}
\begin{exercicios}
  \item $\sin x = \dfrac{1}{2}$
  \item $\cos x = \dfrac{\sqrt{2}}{2}$
  \item $\tan x = 1$
  \item $\sin x = -\dfrac{\sqrt{3}}{2}$
  \item $\cos x = -1$
  \item $\tan x = -\sqrt{3}$
  \item $\sin x = 0$
  \item $\cos x = 2$
  \item $2\cos x + 1 = 0$
  \item $\sqrt{3}\tan x - 1 = 0$
\end{exercicios}

\secao{Equações num Intervalo}{Resolve no intervalo indicado.}
\begin{exercicios}
  \item $\sin x = \dfrac{\sqrt{2}}{2}$, \; em $[0, 2\pi]$
  \item $\cos x = \dfrac{1}{2}$, \; em $[-\pi, \pi]$
  \item $\tan x = \sqrt{3}$, \; em $[0, 2\pi]$
  \item $2\sin x - \sqrt{3} = 0$, \; em $[0, 3\pi]$
  \item $\cos(2x) = 0$, \; em $[0, \pi]$
  \item $\sin\left(x - \dfrac{\pi}{4}\right) = 1$, \; em $[0, 2\pi]$
\end{exercicios}

\secao{Equações do Tipo $\sin u = \sin v$ e $\cos u = \cos v$}{Resolve em $\R$.}
\begin{exercicios}
  \item $\sin(2x) = \sin x$
  \item $\cos(3x) = \cos x$
  \item $\sin\left(x + \dfrac{\pi}{3}\right) = \sin\left(\dfrac{\pi}{6}\right)$
  \item $\cos(2x) = \cos\left(x - \dfrac{\pi}{2}\right)$
  \item $\tan(2x) = \tan\left(x + \dfrac{\pi}{4}\right)$
  \item $\sin x = \cos x$
\end{exercicios}

\secao{Equações que se Reduzem a Equações do 2.º Grau}{Resolve em $\R$.}
\begin{exercicios}
  \item $2\sin^2 x - \sin x = 0$
  \item $\cos^2 x = \dfrac{1}{4}$
  \item $2\cos^2 x - 3\cos x + 1 = 0$
  \item $2\sin^2 x + \sin x - 1 = 0$
  \item $2\cos^2 x + \sin x - 1 = 0$ \; (usa $\cos^2 x = 1 - \sin^2 x$)
  \item $\tan^2 x - 3 = 0$
\end{exercicios}

\secao{Fórmulas de Adição e Duplicação}{Usa $\sin(a \pm b) = \sin a\cos b \pm \cos a \sin b$, $\;\cos(a \pm b) = \cos a\cos b \mp \sin a\sin b$, $\;\sin(2a) = 2\sin a\cos a$ e $\cos(2a) = \cos^2 a - \sin^2 a$.}
\begin{exercicios}
  \item Calcula o valor exato de $\sin 75^\circ$ (usa $75^\circ = 45^\circ + 30^\circ$).
  \item Calcula o valor exato de $\cos 15^\circ$.
  \item Calcula o valor exato de $\cos \dfrac{7\pi}{12}$.
  \item Sabendo que $\sin \alpha = \dfrac{3}{5}$, com $\alpha$ agudo, calcula $\sin(2\alpha)$ e $\cos(2\alpha)$.
  \item Simplifica $\sin x \cos(2x) + \cos x \sin(2x)$.
  \item Mostra que $\cos(2x) = 1 - 2\sin^2 x$.
  \item Resolve: $\sin(2x) = \cos x$.
  \item Resolve: $\cos(2x) + \cos x = 0$.
\end{exercicios}

\secao{Desafios}{Resolve os problemas.}
\begin{exercicios}
  \item Resolve em $[0, 2\pi]$: $\sin x + \cos x = 1$. (Sugestão: eleva ao quadrado e verifica as soluções.)
  \item Resolve em $\R$: $\sqrt{3}\sin x + \cos x = 1$ (escreve o primeiro membro na forma $2\sin(x + \varphi)$).
  \item Mostra que $\tan(a + b) = \dfrac{\tan a + \tan b}{1 - \tan a \tan b}$.
  \item A altura de uma onda é $h(t) = 2\sin\left(\dfrac{\pi t}{3}\right)$, em metros. Em que instantes de $[0, 12]$ a altura é 1\,m?
\end{exercicios}

\end{document}
